\vspace{-5mm}
\chapter{哈巴谷書}
\vspace{-3mm}



\section{神差迦勒底人討罰的启示 1:1-2:1}
\subsection{為公理顛倒呼求,神示迦勒底人必行討罰 1:1-11}
\begin{table}[H]
\centering
\begin{tabular}{p{5.75cm}|p{11.50cm}}\hline\hline
\multicolumn{1}{c|}{哈巴谷為公理顛倒呼求神1-4}&\multicolumn{1}{c}{神示迦勒底人必行討罰5-11}\\\hline
\multirow{2}{5.86cm}{先知哈巴谷所得的默示。2 他說：耶和華啊，我呼求你，你不應允，要到幾時呢？我因強暴哀求你，你還不拯救。 3 你為何使我看見罪孽？你為何看著奸惡而不理呢？毀滅和強暴在我面前，又起了爭端和相鬥的事。 4 因此律法放鬆，公理也不顯明。惡人圍困義人，所以公理顯然顛倒。}
&耶和華說：「你們要向列國中觀看，大大驚奇。因為在你們的時候，我行一件事，雖有人告訴你們，你們總是不信。\\\cline{2-2}
&我必興起迦勒底人，就是那殘忍暴躁之民，通行遍地，占據那不屬自己的住處。 7 他威武可畏，判斷和勢力都任意發出。 8 他的馬比豹更快，比晚上的豺狼更猛。馬兵踴躍爭先，都從遠方而來，他們飛跑如鷹抓食。 9 都為行強暴而來，定住臉面向前，將擄掠的人聚集，多如塵沙。 10 他們譏誚君王，笑話首領，嗤笑一切保障，築壘攻取。 11 他以自己的勢力為神，像風猛然掃過，顯為有罪。」\\\hline
\end{tabular}
\end{table}

\subsection{哀嘆迦勒底人之惡較其所討者尤甚 1:12-17/2:1}
\begin{itemize} 
\itemsep-.25em 
\item 耶和華我的神，我的聖者啊，你不是從亙古而有嗎？我們必不致死。
\item 耶和華啊，你派定他為要刑罰人；磐石啊，你設立他為要懲治人。你眼目清潔，不看邪僻，不看奸惡。行詭詐的，你為何看著不理呢？惡人吞滅比自己公義的，你為何靜默不語呢？你為何使人如海中的魚，又如沒有管轄的爬物呢？  
\item 他用鉤鉤住，用網捕獲，用拉網聚集他們。因此他歡喜快樂， 就向網獻祭，向網燒香，因他由此得肥美的份和富裕的食物。他豈可屢次倒空網羅，將列國的人時常殺戮，毫不顧惜呢？
\item 我要站在守望所，立在望樓上觀看，看耶和華對我說什麼話，我可用什麼話向他訴冤(回答所疑問的)。
\end{itemize}


\section{默示有定日快要應驗 2:2-20}
%\vspace{-5mm}
\subsection{義人因信得生和迦勒底人之禍 2:2-6}
他對我說：「將這默示明明地寫在版上，使讀的人容易讀(隨跑隨讀)。 3 因為這默示有一定的日期，快要應驗，並不虛謊。雖然遲延，還要等候，因為必然臨到，不再遲延。

迦勒底人自高自大，心不正直，\textcolor{blue}{唯義人因信得生}。 5 迦勒底人因酒詭詐、狂傲，不住在家中，擴充心欲好像陰間。他如死不能知足，聚集萬國，堆積萬民，都歸自己。 6 這些國的民豈不都要題起詩歌並俗語，譏刺他說：『\textcolor{red}{禍哉}，迦勒底人！你增添不屬自己的財物，多多取人的當頭，要到幾時為止呢？』 

\subsection{禍哉——為本家積蓄不義之財 2:7-14}
7 咬傷你的豈不忽然起來，擾害你的豈不興起，你就做他們的擄物嗎？ 8 因你搶奪許多的國，殺人流血，向國內的城並城中一切居民施行強暴，所以各國剩下的民都必搶奪你。
為本家積蓄不義之財，在高處搭窩，指望免災的\textcolor{red}{有禍了}！ 10 你圖謀剪除多國的民，犯了罪，使你的家蒙羞，自害己命。 11 牆裡的石頭必呼叫，房內的棟梁必應聲。以人血建城，以罪孽立邑的有禍了！ 13 眾民所勞碌得來的被火焚燒，列國由勞乏而得的歸於虛空，不都是出於萬軍之耶和華嗎？ 14 \textcolor{blue}{認識耶和華榮耀的知識要充滿遍地，好像水充滿洋海一般}。

\subsection{禍哉——為本家積蓄不義之財 2:15-17}
「給人酒喝，又加上毒物使他喝醉，好看見他下體的\textcolor{red}{有禍了}！ 16 你滿受羞辱，不得榮耀。你也喝吧！顯出是未受割禮的。耶和華右手的杯必傳到你那裡，你的榮耀就變為大大的羞辱。 17 你向黎巴嫩行強暴與殘害驚嚇野獸的事，必遮蓋你，因你殺人流血，向國內的城並城中一切居民施行強暴。

\subsection{禍哉——雕刻啞巴偶像 2:18-20}
「雕刻的偶像，人將它刻出來，有什麼益處呢？鑄造的偶像就是虛謊的師傅，製造者倚靠這啞巴偶像有什麼益處呢？ 19 對木偶說『醒起！』，對啞巴石像說『起來！』，那人\textcolor{red}{有禍了}！這個還能教訓人嗎？看哪，是包裹金銀的，其中毫無氣息。  \textcolor{blue}{唯耶和華在他的聖殿中，全地的人都當在他面前肅敬靜默}。」

\section{哈巴谷因明白默示而祈禱 3:1-19}

先知哈巴谷的禱告，調用流離歌。

2 耶和華啊，我\textcolor{red}{聽見}你的名聲（言語）就懼怕。耶和華啊，求你在這些年間復興你的作為，在這些年間顯明出來，在發怒的時候以憐憫為念。 3 神從提幔而來，聖者從巴蘭山臨到。（細拉）他的榮光遮蔽諸天，頌讚充滿大地。 4 他的輝煌如同日光，從他手裡射出光線，在其中藏著他的能力。 5 在他前面有瘟疫流行，在他腳下有熱症發出。 6 他站立，量了大地（使地震動）；觀看，趕散萬民。永久的山崩裂，長存的嶺塌陷，他的作為與古時一樣。 
7 我\textcolor{red}{見}古珊的帳篷遭難，米甸的幔子戰兢。 8 耶和華啊，你乘在馬上，坐在得勝的車上，豈是不喜悅江河，向江河發怒氣，向洋海發憤恨嗎？ 9 你的弓全然顯露，向眾支派所起的誓都是可信的。（細拉）你以江河分開大地。 10 山嶺\textcolor{red}{見}你，無不戰懼；大水氾濫過去，深淵發聲，洶湧翻騰（向上舉手）。 11 因你的箭射出發光，你的槍閃出光耀，日月都在本宮停住。 12 你發憤恨通行大地，發怒氣責打列國如同打糧。 13 你出來要拯救你的百姓，拯救你的受膏者，打破惡人家長的頭，露出他的腳 （根基），直到頸項。（細拉） 14 你用敵人的戈矛刺透他戰士的頭，他們來如旋風，要將我們分散，他們所喜愛的是暗中吞吃貧民。 15 你乘馬踐踏紅海，就是踐踏洶湧的大水。

我\textcolor{red}{聽見}耶和華的聲音，身體戰兢，嘴唇發顫，骨中朽爛，我在所立之處戰兢。我只可安靜等候災難之日臨到，犯境之民上來。 17 雖然無花果樹不發旺，葡萄樹不結果，橄欖樹也不效力，田地不出糧食，圈中絕了羊，棚內也沒有牛， 18 然而我要因耶和華歡欣，因救我的神喜樂。 19 主耶和華是我的力量，他使我的腳快如母鹿的蹄，又使我穩行在高處。

這歌交於伶長，用絲弦的樂器。
